\documentclass[11pt,a4paper]{article}
\PassOptionsToPackage{pdfusetitle=false}{hyperref}
\usepackage[utf8]{inputenc}
\usepackage[T1]{fontenc}
\usepackage[margin=2.5cm]{geometry}
\usepackage{amsmath,amssymb}
\usepackage{booktabs,array}
\usepackage{xcolor}
\usepackage{tcolorbox}
\tcbuselibrary{skins,breakable}
\usepackage{enumitem}
\usepackage[pdfusetitle=false]{hyperref}
\usepackage[numbers]{natbib}
\bibliographystyle{unsrtnat}
\hypersetup{colorlinks=true,linkcolor=blue!60!black,urlcolor=blue!60!black,
            citecolor=blue!60!black,
            pdftitle={The GTE--Wolfram Bridge},
            pdfauthor={Nova Spivack}}
\setlength{\emergencystretch}{2em}
\usepackage{tikz}
\usetikzlibrary{arrows.meta,positioning,shapes.geometric,calc,fit,backgrounds}

%─── Color palette ───────────────────────────────────────────────────────────
\definecolor{boxblue}{RGB}{220,235,255}
\definecolor{boxorange}{RGB}{255,240,210}
\definecolor{boxgreen}{RGB}{220,245,220}
\definecolor{boxpurple}{RGB}{240,225,255}
\definecolor{boxtan}{RGB}{255,245,220}
\definecolor{boxteal}{RGB}{210,240,238}

%─── Macros ──────────────────────────────────────────────────────────────────
\newcommand{\term}[1]{\textbf{#1}}
\newcommand{\lean}[1]{%
  \penalty-100
  \begingroup\def\_{\textunderscore\allowbreak}\texttt{#1}\endgroup}
\newcommand{\CatAL}{\textbf{CatAL}}
\newcommand{\CatAD}{\textbf{CatAD}}
\newcommand{\CatA}{\textbf{CatA}}
\newcommand{\PHIMDL}{$\Phi_{\rm MDL}$}
\newcommand{\PhiMDL}{\Phi_{\rm MDL}}
\newcommand{\poly}{p}

\newcommand{\TutorialDOI}{10.5281/zenodo.20669164}
\date{2026}

\title{\textbf{The GTE--Wolfram Bridge}\\[6pt]
\large Rule~110, the Causal Graph, and the Selection Principle\\[8pt]
\normalsize
\href{https://novaspivack.com/research}{novaspivack.com/research}\texorpdfstring{\quad}{ }
\href{https://doi.org/10.5281/zenodo.20560550}{P48 complete theory monograph}\\[4pt]
\small\href{https://doi.org/\TutorialDOI}{doi.org/\TutorialDOI}
  \textit{(registration pending)}\\[4pt]
\normalsize\textit{UGP Physics Tutorial Series}}
\author{Nova Spivack}

\begin{document}
\setcounter{tocdepth}{2}
\maketitle
\tableofcontents
\newpage

\begin{tcolorbox}[colback=blue!6,colframe=blue!40!black,
  title=\textbf{UGP Physics Tutorial Series}]
\textbf{Prerequisites (read first):}
  \href{https://doi.org/10.5281/zenodo.20657889}%
    {\emph{The GTE Polynomial: A Step-by-Step Cheat Sheet}}
  $\cdot$
  \href{https://doi.org/10.5281/zenodo.20657919}%
    {\emph{How the UWCA Works}}
  $\cdot$
  \href{https://doi.org/10.5281/zenodo.20657906}%
    {\emph{The Three-Tape CMCA}}\\[4pt]
\textbf{Next tutorials:}
  \href{https://doi.org/10.5281/zenodo.20669142}%
    {\emph{Falsifiability: How to Test the GTE Framework}}\\[4pt]
\textbf{See also:}
  \href{https://doi.org/10.5281/zenodo.20657872}%
    {\emph{The Levels of the Theory}}
  $\cdot$
  \href{https://doi.org/10.5281/zenodo.20657876}%
    {\emph{How MDL Selects the 19-Bit Polynomial}}
  $\cdot$
  \href{https://doi.org/10.5281/zenodo.20657895}%
    {\emph{Perfect Self-Containment and MDL}}
\end{tcolorbox}

\begin{tcolorbox}[colback=boxgreen,colframe=green!50!black,
  title=\textbf{Claim-Strength Taxonomy}]
Throughout this tutorial, results are labeled with a confidence grade:
\begin{itemize}[itemsep=2pt]
  \item \textbf{$\mathbf{A}_{\rm Lean}$ (CatAL)} --- Machine-certified in Lean~4 with zero \texttt{sorry} and zero custom axioms.  Independently verifiable by anyone with Lean installed.
  \item \textbf{A/D (CatAD)} --- Analytically derived with full derivation, computationally confirmed.  Not yet machine-certified.
  \item \textbf{A (CatA)} --- Analytically derived; derivation complete but not independently machine-certified.
  \item \textbf{A/D (conditional)} --- Derived under a named, stated premise; the premise is itself an open problem or a CatB assumption.
  \item \textbf{B (CatB)} --- A stated working hypothesis or structural premise; not yet derived.
\end{itemize}
\end{tcolorbox}

%══════════════════════════════════════════════════════════════════════════════
\section{The Wolfram Physics Programme}
\label{sec:wolfram_programme}
%══════════════════════════════════════════════════════════════════════════════

\begin{tcolorbox}[colback=boxblue,colframe=blue!40!black,
  title=What this tutorial addresses]
The Wolfram Physics Project seeks a fundamental computational rule that generates
spacetime and matter.  GTE provides what Wolfram Physics lacks: a \emph{selection
principle} that identifies the unique such rule from first principles, and
machine-certifies it.  This tutorial explains the bridge between the two programmes.
\end{tcolorbox}

\subsection{The Wolfram Physics Approach}

The \term{Wolfram Physics Project} (Stephen Wolfram and Jonathan Gorard, 2020) is
the most recent systematic attempt to derive physics from a discrete computational
rule.  Its three key components are:

\begin{enumerate}[itemsep=4pt]
  \item \textbf{The hypergraph rule.}  The universe is modeled as a hypergraph
    (a collection of points with multi-way connections) that evolves by a local
    rewriting rule.  Each application of the rule creates new connections and
    deletes old ones.

  \item \textbf{The causal graph.}  The \term{causal graph} records which rule
    applications causally depend on which others — essentially, which events
    produce the inputs for which later events.  In the continuum limit, the
    causal graph encodes a Lorentzian metric on spacetime.

  \item \textbf{Multiway systems.}  When multiple rule applications are possible,
    the \term{multiway system} tracks all of them simultaneously.  Quantum
    mechanics is proposed to emerge from the structure of paths through the
    multiway graph.
\end{enumerate}

The programme has produced a consistent framework for how spacetime geometry and
general relativity can emerge from a discrete computational rule.  Gorard and
Wolfram have shown that Einstein's equations appear in the continuum limit of
appropriate causal graphs \cite{SpivackGTECompleteFramework}.

\subsection{The Central Gap: Which Rule?}

The Wolfram Physics programme is maximally inclusive: it defines the
\term{Ruliad} as the union of all possible computations.  The physical universe
is somewhere in the Ruliad; the question is \emph{where} — which specific rule?

The Wolfram approach provides no answer to this.  The universe is not in the
Ruliad because of some meta-selection process; it simply is where it is.
This leaves the selection question open, which is the same open question that
all of digital physics has faced since Konrad Zuse first proposed computational
physics in 1969.

GTE addresses this gap directly.

%══════════════════════════════════════════════════════════════════════════════
\section{What GTE Adds: The Selection Principle}
\label{sec:selection}
%══════════════════════════════════════════════════════════════════════════════

\begin{tcolorbox}[colback=boxorange,colframe=orange!60!black,
  title=GTE's answer to ``which rule?'']
The unique MDL-minimal computational rule consistent with the PSC constraints
is $p(L,C,R) = C + R - CR - LCR$ over $\mathrm{GF}(7)$.  This rule is selected
not by the observer's location in a Ruliad, but by an objective mathematical
criterion: Minimum Description Length.
\end{tcolorbox}

\term{GTE} (\term{Generative Triple Evolution}) applies two
interacting selection criteria:

\textbf{PSC} (\term{Perfect Self-Containment}):  The candidate rule must be
the unique rule that can self-referentially justify its own selection from
all possible rules.  PSC is a set of five axioms about which theories can be
considered complete without external input.  It eliminates all but a vanishing
fraction of possible rules.

\textbf{MDL} (\term{Minimum Description Length}):  Among all PSC-admissible
rules, the one with the shortest complete description is selected.  MDL is the
operational expression of \term{PI} (Presentation Invariance), the fifth PSC
axiom.

Together, PSC and MDL uniquely identify the GTE polynomial $p(L,C,R)$ over
$\mathrm{GF}(7)$ as the rule governing our universe.  This is not a heuristic
argument — it is a derivation chain, with the key steps certified in Lean~4.

\textbf{What this gives Wolfram Physics:} A derivation, not a search.  Instead
of exploring the space of rules and comparing to observation, GTE proves which
rule is unique and then verifies it matches observation.

%══════════════════════════════════════════════════════════════════════════════
\section{The Triangle: Two Paths to Rule 110}
\label{sec:triangle}
%══════════════════════════════════════════════════════════════════════════════

The structure of GTE contains a remarkable internal triangle connecting three
objects through two independent paths, both converging on Rule~110.

\subsection{The Two Objects}

\begin{itemize}[itemsep=4pt]
  \item \textbf{The UGP arithmetic cascade.}  The GTE framework starts from the
    UGP (\term{Universe Generating Polynomial}) arithmetic — a specific sequence
    of integer triples generated by the cascade map $T$, starting from the
    Lepton Seed $(1,73,823)$.  This cascade is derived from the
    $N_{\rm gen} = 3$ and $N_{\rm fam} = 5$ generation/family counts.

  \item \textbf{The GTE polynomial $p$.}  The polynomial
    $p(L,C,R) = C + R - CR - LCR$ over $\mathrm{GF}(7)$ — the algebraic
    certificate of the GTE framework.
\end{itemize}

These look like independent objects.  But they are not: there are two distinct
paths from UGP to $p$, both ending at exactly Rule~110.

\subsection{Path A: Direct-Interpolation Lift}

\begin{tcolorbox}[colback=boxteal,colframe=blue!40!black,
  title={Direct-Interpolation Lift Theorem (\CatAL)}]
\textbf{Statement:} The total-parity shadow of the canonical cascade triples
constrains all 8 binary input-output pairs of the GTE polynomial.
Combined with MDL minimality (which enforces the vacuum-transparency condition),
the unique multilinear rule consistent with these constraints is exactly
$p(L,C,R) = C + R - CR - LCR$ over $\mathrm{GF}(7)$.  Rule~110 is a corollary
(the binary restriction of $p$).\\[4pt]
Machine-certified: LT-088-65--69 (\CatAL, zero sorry, zero custom axioms).
\end{tcolorbox}

The derivation in four steps:

\textbf{Step 1 — The parity shadow.}  Reduce every coordinate of every cascade
triple modulo 2 — keep only whether each value is even (0) or odd (1).
This produces a sequence of binary triples.  Collecting the unique ones from the
ten canonical cascade triples gives 7 distinct binary $(L,C,R)$ patterns,
each paired with an output bit also read from the cascade.

\textbf{Step 2 — The 8th row.}  One binary input pattern is not covered by the
parity shadow: $(1,1,1)$ (all three neighbors odd).  The vacuum-transparency
condition — $p(0,0,0) = 0$ (the all-zeros state is a fixed point) — forces this
remaining output to be 0.  Now all 8 binary rows are specified.

\textbf{Step 3 — Polynomial interpolation.}  A multilinear polynomial of
degree $\leq 3$ over three binary variables has exactly $2^3 = 8$ basis
monomials: $\{1, L, C, R, LC, LR, CR, LCR\}$.  With 8 known input-output
pairs and 8 unknowns, the linear system
$\mathbf{V}\,\mathbf{c} = \mathbf{y}$ over $\mathrm{GF}(7)$ has a unique
solution (the Vandermonde-type matrix is full rank).

\textbf{Step 4 — The result.}  The unique solution is exactly
$p(L,C,R) = C + R - CR - LCR$ over $\mathrm{GF}(7)$,
and the binary restriction $p|_{\{0,1\}^3}$ produces the Rule~110 truth table.

\textbf{Worked numerical check:}
For input $(L,C,R) = (0,1,1)$: $p = 1+1-1\cdot1-0\cdot1\cdot1 = 1$.
Rule~110 for input $011_2 = 3$: output bit is $1$. $\checkmark$

\subsection{Path B: MDL Direct Selection}

Path B is the standard MDL derivation.  Among all $k=7$, $r=1$ cellular
automaton rules consistent with the PSC constraints (Standard Model
derivability criterion), the shortest description is the GTE polynomial,
requiring exactly 19 bits:

\begin{itemize}[itemsep=2pt]
  \item 8 bits — the Rule~110 truth table (binary restriction)
  \item 3 bits — the prime modulus ($\lceil\log_2 7\rceil = 3$)
  \item 5 bits — the monomial selector (which 4 of 8 basis terms are nonzero)
  \item 1 bit — chirality
  \item 2 bits — the $-1$ coefficient signs
\end{itemize}

A naive lookup table requires $343 \times \lceil\log_2 7\rceil \approx 963$
bits — about 50 times larger.  The MDL Uniqueness Theorem proves no shorter
valid specification exists (\CatAL) \cite{SpivackGTEPolynomialWolfram}.

\subsection{Convergence: Both Paths Identify Rule 110}

The Direct-Interpolation Lift (Path A) shows that the UGP orbit's arithmetic
structure forces the Rule~110 truth table on the polynomial's 8 binary rows.
MDL selection (Path B) then forces the unique multilinear extension to all
343 GF(7) inputs.  \emph{Rule~110 is a theorem of GTE, not a postulate.}

This is the triangle: UGP arithmetic $\to$ orbit parity shadow $\to$ 8 binary
rows $\to$ (by MDL interpolation) $\to$ unique polynomial $p$ $\to$ (binary
restriction) $\to$ Rule~110.  The path is sequential and fully derived.

%══════════════════════════════════════════════════════════════════════════════
\section{The Chirality Census: \{Rule 110, Rule 124\}}
\label{sec:chirality}
%══════════════════════════════════════════════════════════════════════════════

\begin{tcolorbox}[colback=boxpurple,colframe=blue!40!black,
  title={Chirality Census Theorem (\CatAL)}]
Over all 120 orderings of the three SM fermion generations, the survivor set
of rules consistent with the GTE orbit parity shadow is exactly
$\{\text{Rule~110},\; \text{Rule~124}\}$ — the chiral pair related by
left-right reflection.  The choice of generation ordering convention is the
\emph{chirality gauge}: it selects one of the two chiral orientations but does
not change any physical observable.
\end{tcolorbox}

Rule~110 and Rule~124 are related by the AGL(1,7) reflection symmetry of the
$\mathrm{GF}(7)$ alphabet (the map $x \mapsto -x \pmod{7}$).  Machine-certified:
\lean{orbit\_chirality\_census} (\CatAL) \cite{SpivackGTECompleteFramework}.

\textbf{Why this matters for Wolfram Physics:} The Wolfram Physics programme
works with causal graphs that are directional.  The left-right orientation of
a CA rule is precisely the chirality that GTE identifies as a gauge degree of
freedom.  The fact that GTE derives exactly two chiral orientations (not one,
not many) connects naturally to the chiral structure of the Standard Model's
weak interactions.

\textbf{Worked example.}  Rule~110 has the bit pattern $01101110_2 = 110_{10}$
(reading binary outputs from $(1,1,1)$ to $(0,0,0)$ in order).  Rule~124 has
the bit pattern $01111100_2 = 124_{10}$.  One is the left-right reflection of
the other over the $\{0,1\}$ alphabet.  Both are derived from the UGP orbit;
neither is independently postulated.

%══════════════════════════════════════════════════════════════════════════════
\section{The Causal Graph: 1023 Nodes and Queen Anne's Lace}
\label{sec:causal}
%══════════════════════════════════════════════════════════════════════════════

\subsection{Generating the Causal Graph}

The Wolfram Physics Project provides a framework for computing the causal graph
of a hypergraph rule.  For GTE, the relevant rule is \texttt{ruleGTE}: a
WolframModel hyperedge rule encoding two overlapping CA triplet neighborhoods:
\begin{equation}
  \{a,b,c\},\{c,d,e\} \;\to\; \{a,b,f\},\{f,c,d\},\{e,b,g\},\{g,d,h\}.
  \label{eq:rulegte}
\end{equation}

Running WolframModel for 10 generations produces a causal graph with
\textbf{1023 nodes and 2044 edges} — a perfect binary tree of depth 10
\cite{SpivackGTECompleteFramework}.  Machine-certified:
\begin{itemize}[itemsep=2pt]
  \item \lean{perfectTree\_numNodes}: $2^n - 1$ nodes by induction (\CatAL).
  \item \lean{gte\_rulegte\_ten\_generations}: 1023 events at depth 10 (\CatAL).
  \item \lean{perfectTree\_horton\_ratio}: Horton ratio $r_B = 2$ at all levels (\CatAL).
\end{itemize}

\subsection{The Queen Anne's Lace Structural Equivalence}

The GTE causal graph (10 generations, 1023 nodes) is structurally equivalent to
the Lindenmayer system $F \to F[+F][-F]$ at $\theta \approx 25^\circ$, $r \approx 0.70$
— the canonical L-system model of \emph{Daucus carota} (Queen Anne's lace, a
flowering plant in the Apiaceae family) \cite{SpivackGTECompleteFramework}.

Five independent metrics agree:
\begin{itemize}[itemsep=2pt]
  \item Horton branching ratio $r_B = 2$ (\CatAL)
  \item Strahler number 10
  \item Fractal dimension $D \approx 1.943$
  \item Depth 10
  \item Bilateral symmetry
\end{itemize}

The shared computational primitive: \texttt{ruleGTE} creates 2 overlapping pairs
from 1 (binary branching), exactly as meristematic apical cell division creates
2 daughter cells.  This is a structural equivalence (\CatA) — a shared graph
topology — not a formal isomorphism between cellular automaton dynamics and
plant growth.  The physical significance is that the GTE causal graph's fractal
structure has the Horton ratio $r_B = 2$ machine-certified as a theorem of the
rule's combinatorics, not as a measured fit.

%══════════════════════════════════════════════════════════════════════════════
\section{Three Key Distinctions for Wolfram Physics Readers}
\label{sec:distinctions}
%══════════════════════════════════════════════════════════════════════════════

Readers familiar with the Wolfram Physics Project should note three precise
differences between GTE and the Wolfram approach.

\subsection{Distinction 1: The CA Is the Certificate, Not the Substrate}

The most important difference.  The Wolfram Physics Project proposes that the
universe \emph{is} a hypergraph rule evolving on a hypergraph.

GTE proves this cannot be exactly right: no finite cellular automaton can be the
exact physical substrate of our universe.\\
Machine-certified: \lean{no\_finite\_ca\_exact\_lorentz\_replica}
(\CatAL, zero sorry) \cite{SpivackCompUniversality}.

The reason: every CA with a finite lattice spacing has Lorentz-violating
artifacts at that spacing.  The physical substrate must recover exact Lorentz
invariance, which no lattice CA achieves.

The GTE resolution is the two-level architecture:
\begin{itemize}[itemsep=3pt]
  \item \textbf{Level 1 (certificate):} The CMCA (the discrete CA) encodes the
    structural constraints of the physical field via the Algebraic Lifting
    Theorem (\CatAL).  It is an algebraic certificate, not the substrate.
  \item \textbf{Level 2 (substrate):} The $\Phi_{\rm MDL}$ field is the actual
    physical substrate — a continuous $\mathbb{Z}_7$-symmetric Klein-Gordon
    field on $\mathbb{R}^{3+1}$.
\end{itemize}
GTE does not claim the universe IS a CA.  It proves the universe's physics has an
algebraic certificate that IS a CA.  These are different claims, and the
distinction matters enormously for Lorentz invariance and quantum mechanics.

\subsection{Distinction 2: MDL Uniquely Selects the GTE Polynomial}

The Wolfram Physics Project explores the space of rules computationally and
compares to observation.  GTE selects the rule analytically.

The polynomial $\poly = C + R - CR - LCR$ over $\mathrm{GF}(7)$ is the unique
MDL-minimal description consistent with the PSC constraints — machine-certified
in Lean~4, zero sorry \cite{SpivackGTEPolynomialWolfram}.  Every
$\mathbb{Z}_N \times \mathbb{Z}_M$ competitor has been enumerated and
eliminated.  MDL is the selection principle that digital physics has lacked.

This means GTE does not search for rules; it derives the rule.

\subsection{Distinction 3: Rule 110 Is a Theorem, Not a Postulate}

In Wolfram's programme, Rule~110 was discovered by computational enumeration and
its Turing universality proved by Matthew Cook (1994, published 2004).  It is an
empirical discovery about which simple CA is computationally interesting.

In GTE, Rule~110 is a theorem:
\begin{equation*}
  p\big|_{\{0,1\}^3} = \text{Rule~110},
\end{equation*}
machine-certified: \lean{rule110\_z7\_poly\_rep} (\CatAL, zero sorry,
\texttt{native\_decide}) \cite{SpivackGTEPolynomialWolfram}.  Rule~110 is the
binary restriction of the GTE polynomial, derived from the UGP orbit arithmetic.
Cook's proof of Rule~110's Turing universality follows as a structural consequence
(\lean{phimdl\_turing\_universal}, \CatAL)
\cite{SpivackCompUniversality}.

%══════════════════════════════════════════════════════════════════════════════
\section{The 19-Bit Program: Connecting MDL to Wolfram Language}
\label{sec:19bit}
%══════════════════════════════════════════════════════════════════════════════

A central claim of GTE is that the universe's complete algorithmic description
requires \textbf{19 bits}.  This is not an approximation; it is the MDL complexity
of the GTE polynomial in its natural description class.

\subsection{The 19-Bit Decomposition}

\begin{table}[htbp]
\centering
\caption{The 19-bit decomposition of the GTE polynomial description.
Each bit count is exact and follows from the structure of the polynomial
representation class.}
\label{tab:19bits}
\begin{tabular}{lll}
\toprule
Component & Bits & What it encodes \\
\midrule
Rule~110 truth table & 8 & Binary outputs for all 8 binary inputs \\
Prime modulus & 3 & $\lceil\log_2 7\rceil$ — selects the GF(7) alphabet \\
Monomial selector & 5 & Which 4 of 8 basis monomials are nonzero \\
Chirality & 1 & Left-right orientation (the chirality gauge) \\
Coefficient signs & 2 & The $-1$ coefficients on $CR$ and $LCR$ \\
\midrule
\textbf{Total} & \textbf{19} & The complete GTE polynomial \\
\bottomrule
\end{tabular}
\end{table}

For comparison, specifying the rule as a lookup table requires
$343 \times \lceil\log_2 7\rceil = 343 \times 3 \approx 963$ bits.
The polynomial representation is $963/19 \approx 50\times$ more compressed.

\subsection{Wolfram Language Implementation}

The three-tape CMCA that implements the GTE polynomial can be written directly
in Wolfram Language using \texttt{CellularAutomaton} and \texttt{WolframModel}.
The rule \texttt{ruleGTE} produces the causal graph discussed in
\S\ref{sec:causal}.  The polynomial's 343 outputs can be verified by:

\begin{center}
\texttt{Table[Mod[C + R - C*R - L*C*R, 7], \{L,0,6\}, \{C,0,6\}, \{R,0,6\}]}
\end{center}

This generates the complete GTE update table from the polynomial formula.  The
Rule~110 connection is verified by restricting to $\{L,C,R\} \in \{0,1\}^3$:
the resulting 8 outputs match the Rule~110 truth table exactly.

%══════════════════════════════════════════════════════════════════════════════
\section{Open Questions at the Bridge}
\label{sec:open}
%══════════════════════════════════════════════════════════════════════════════

Several questions at the GTE--Wolfram interface remain open.  Presenting them
honestly is essential to the scientific record.

\subsection{The Causal-Graph / Gorard Chain Correspondence}

The Wolfram Physics programme identifies the emergent spacetime metric with the
limiting structure of the causal graph.  Gorard's result shows that Einstein's
equations appear in the continuum limit of appropriate causal graphs.

GTE uses a related object: the \term{Gorard chain} (defined in P48) is a
sequence of causal events in the $\Phi_{\rm MDL}$ field.  The \emph{precise
correspondence} between the P48 Gorard-chain limit and the Wolfram Physics
causal graph is stated as an open question in the GTE framework.  Both
directions — GTE's and Wolfram's — point toward the same Einstein equations;
whether they can be formally identified as the same limiting structure is not
yet proven \cite{SpivackGTECompleteFramework}.

\subsection{Quantum Mechanics from the Multiway System}

The Wolfram Physics programme proposes that quantum mechanics emerges from paths
through the multiway evolution graph.  GTE derives quantum mechanics — including
the Born rule (\CatAL, \lean{born\_rule\_unconditional}) — from the
$[D]$-weighted path integral over kink configurations.

Whether the GTE $[D]$-weighting is the same mathematical object as the
Wolfram-Gorard branch weights in the multiway graph is a precise open question
at the interface of the two frameworks.

%══════════════════════════════════════════════════════════════════════════════
\section{Summary: What GTE Offers Wolfram Physics}
\label{sec:summary}
%══════════════════════════════════════════════════════════════════════════════

\begin{tcolorbox}[colback=boxgreen,colframe=green!50!black,
  title=The five contributions]
\begin{enumerate}[itemsep=4pt]
  \item \textbf{A derivation, not a search.}  GTE identifies the unique rule
    (the GTE polynomial over GF(7)) from first principles, via PSC and MDL.
    No exploration of rule space is required.

  \item \textbf{Machine certification.}  The key steps — MDL uniqueness, Rule~110
    as a theorem, the chirality census — are Lean~4 certified, zero sorry.
    The derivation chain is formally verified.

  \item \textbf{A 19-bit description.}  The GTE polynomial is the minimum
    description of the universe's computational rule.  This connects Wolfram's
    emphasis on simple rules to a precise bit-count theorem.

  \item \textbf{The CA as certificate, not substrate.}  GTE resolves the
    Lorentz-invariance tension in discrete physics by separating the algebraic
    certificate (the CMCA) from the physical substrate (the $\Phi_{\rm MDL}$
    field).

  \item \textbf{A causal graph with certified structure.}  The GTE causal graph
    from \texttt{ruleGTE} has machine-certified tree structure, Horton ratio, and
    an explicit structural equivalence to a botanical branching system.
\end{enumerate}
\end{tcolorbox}

\section*{Further Reading}
\begin{tcolorbox}[colback=orange!8,colframe=orange!50!black,
  title=\textbf{Primary Sources}]
The results in this tutorial are derived in the following papers,
all available open-access on Zenodo and at
\href{https://novaspivack.com/research}{novaspivack.com/research}:
\begin{itemize}[itemsep=2pt]
  \item \textbf{P48 --- The Complete GTE Framework} (main synthesis monograph,
    including the Wolfram Physics context chapter):\\
        \href{https://doi.org/10.5281/zenodo.20560550}{doi.org/10.5281/zenodo.20560550}
  \item \textbf{P49 --- MDL Selects the Wolfram Rule: $\mathbb{Z}_7$ Dynamics}
        (the GTE--Wolfram bridge and the Direct-Interpolation Lift Theorem):\\
        \href{https://doi.org/10.5281/zenodo.20572656}{doi.org/10.5281/zenodo.20572656}
  \item \textbf{P28 --- Computational Universality and the Standard Model}
        (no-CA-replica theorem; Turing universality from GTE):\\
        \href{https://doi.org/10.5281/zenodo.20259513}{doi.org/10.5281/zenodo.20259513}
  \item \textbf{P45 --- The Three-Tape Chiral Minkowski Cellular Automaton}
        (the CMCA as the Level-1 algebraic certificate):\\
        \href{https://doi.org/10.5281/zenodo.20465805}{doi.org/10.5281/zenodo.20465805}
\end{itemize}
Lean~4 machine proofs:
\href{https://github.com/novaspivack/ugp-lean}{github.com/novaspivack/ugp-lean}
\end{tcolorbox}

\bibliography{../../papers/bib/Spivack_Papers_Bibliography}
\end{document}
